
\documentclass{article}
\usepackage{colortbl}
\usepackage{makecell}
\usepackage{multirow}
\usepackage{supertabular}

\begin{document}

\newcounter{utterance}

\centering \large Interaction Transcript for game `wordle', experiment `medium\_frequency\_words\_no\_clue\_no\_critic', episode 4 with claude{-}opus{-}4{-}8{-}azure.
\vspace{24pt}

{ \footnotesize  \setcounter{utterance}{1}
\setlength{\tabcolsep}{0pt}
\begin{supertabular}{c@{$\;$}|p{.15\linewidth}@{}p{.15\linewidth}p{.15\linewidth}p{.15\linewidth}p{.15\linewidth}p{.15\linewidth}}
    \# & \multicolumn{2}{c}{Player} && \multicolumn{2}{c}{Game Master} \\
    \hline

    \theutterance \stepcounter{utterance}  
    & & & \multicolumn{4}{p{0.6\linewidth}}{
        \cellcolor[rgb]{0.9,0.9,0.9}{
            \makecell[{{p{\linewidth}}}]{
                \texttt{\tiny{[P1$\langle$GM]}}
                \texttt{Sei un mago del linguaggio a cui piace indovinare parole usando le regole fornite.} \\
\\ 
\texttt{Benvenuto a Wordle! Hai sei tentativi per indovinare la parola obiettivo, una parola valida di cinque lettere minuscole. Per favore usa i tag "spiegazione:" e "tentativo:" per fornire una spiegazione concisa per ogni tentativo.} \\
\\ 
\texttt{Per esempio, se il tuo tentativo è "parte", la tua risposta dovrebbe essere} \\
\texttt{spiegazione: questa è una comune parola di cinque lettere, e sto iniziando il mio tentativo con questa parola.} \\
\texttt{tentativo: parte} \\
\\ 
\texttt{Dopo ogni tentativo, la tua risposta verrà convalidata e riceverai un feedback che indica quali lettere sono corrette (verde), quali lettere sono corrette ma nella posizione sbagliata (giallo) e quali lettere sono errate (rosso). Questo feedback può essere utile per determinare quali lettere includere o escludere nel tuo prossimo tentativo.} \\
\\ 
\texttt{Per esempio, il feedback per "parte" potrebbe essere:} \\
\texttt{feedback\_tentativo: p<giallo> a<giallo> r<verde> t<giallo> e<rosso>} \\
\\ 
\texttt{La spiegazione dovrebbe contenere dettagli su come il feedback\_tentativo viene usato per arrivare a un nuovo tentativo.} \\
\\ 
\texttt{Non generare alcun altro testo.} \\
\\ 
\texttt{Iniziamo con il tuo primo tentativo.} \\
            }
        }
    }
     \\ \\

    \theutterance \stepcounter{utterance}  
    & \multicolumn{4}{p{0.6\linewidth}}{
        \cellcolor[rgb]{0.561,0.737,0.561}{
            \makecell[{{p{\linewidth}}}]{
                \texttt{\tiny{[P1$\rangle$GM]}}
                \texttt{spiegazione: inizio con una parola comune di cinque lettere che contiene molte vocali per raccogliere informazioni utili.} \\
\texttt{tentativo: carte} \\
            }
        }
    }
    & & \\ \\

    \theutterance \stepcounter{utterance}  
    & & & \multicolumn{2}{p{0.3\linewidth}}{
        \cellcolor[rgb]{0.9,0.9,0.9}{
            \makecell[{{p{\linewidth}}}]{
                \texttt{\tiny{[GM$|$GM]}}
                \texttt{attempts = 1} \\
\texttt{target\_word = nonno} \\
\texttt{guess = carte} \\
\texttt{guess\_feedback = c<red> a<red> r<red> t<red> e<red>} \\
            }
        }
    }
    & & \\ \\

    \theutterance \stepcounter{utterance}  
    & & & \multicolumn{4}{p{0.6\linewidth}}{
        \cellcolor[rgb]{0.9,0.9,0.9}{
            \makecell[{{p{\linewidth}}}]{
                \texttt{\tiny{[P1$\langle$GM]}}
                \texttt{feedback\_tentativo: c<rosso> a<rosso> r<rosso> t<rosso> e<rosso>} \\
\\ 
\texttt{Fornisci la tua risposta solo in questo formato.} \\
\texttt{spiegazione: dettagli} \\
\texttt{tentativo: parola} \\
            }
        }
    }
     \\ \\

    \theutterance \stepcounter{utterance}  
    & \multicolumn{4}{p{0.6\linewidth}}{
        \cellcolor[rgb]{0.561,0.737,0.561}{
            \makecell[{{p{\linewidth}}}]{
                \texttt{\tiny{[P1$\rangle$GM]}}
                \texttt{spiegazione: escludo c, a, r, t, e. Provo una parola con vocali diverse come o, u e nuove consonanti.} \\
\texttt{tentativo: lungo} \\
            }
        }
    }
    & & \\ \\

    \theutterance \stepcounter{utterance}  
    & & & \multicolumn{2}{p{0.3\linewidth}}{
        \cellcolor[rgb]{0.9,0.9,0.9}{
            \makecell[{{p{\linewidth}}}]{
                \texttt{\tiny{[GM$|$GM]}}
                \texttt{attempts = 2} \\
\texttt{target\_word = nonno} \\
\texttt{guess = lungo} \\
\texttt{guess\_feedback = l<red> u<red> n<green> g<red> o<green>} \\
            }
        }
    }
    & & \\ \\

    \theutterance \stepcounter{utterance}  
    & & & \multicolumn{4}{p{0.6\linewidth}}{
        \cellcolor[rgb]{0.9,0.9,0.9}{
            \makecell[{{p{\linewidth}}}]{
                \texttt{\tiny{[P1$\langle$GM]}}
                \texttt{feedback\_tentativo: l<rosso> u<rosso> n<verde> g<rosso> o<verde>} \\
\\ 
\texttt{Fornisci la tua risposta solo in questo formato.} \\
\texttt{spiegazione: dettagli} \\
\texttt{tentativo: parola} \\
            }
        }
    }
     \\ \\

    \theutterance \stepcounter{utterance}  
    & \multicolumn{4}{p{0.6\linewidth}}{
        \cellcolor[rgb]{0.561,0.737,0.561}{
            \makecell[{{p{\linewidth}}}]{
                \texttt{\tiny{[P1$\rangle$GM]}}
                \texttt{spiegazione: n in posizione 3, o in posizione 5. Escludo c,a,r,t,e,l,u,g. Provo "sonno" che ha n in terza e o in quinta posizione.} \\
\texttt{tentativo: sonno} \\
            }
        }
    }
    & & \\ \\

    \theutterance \stepcounter{utterance}  
    & & & \multicolumn{2}{p{0.3\linewidth}}{
        \cellcolor[rgb]{0.9,0.9,0.9}{
            \makecell[{{p{\linewidth}}}]{
                \texttt{\tiny{[GM$|$GM]}}
                \texttt{attempts = 3} \\
\texttt{target\_word = nonno} \\
\texttt{guess = sonno} \\
\texttt{guess\_feedback = s<red> o<green> n<green> n<green> o<green>} \\
            }
        }
    }
    & & \\ \\

    \theutterance \stepcounter{utterance}  
    & & & \multicolumn{4}{p{0.6\linewidth}}{
        \cellcolor[rgb]{0.9,0.9,0.9}{
            \makecell[{{p{\linewidth}}}]{
                \texttt{\tiny{[P1$\langle$GM]}}
                \texttt{feedback\_tentativo: s<rosso> o<verde> n<verde> n<verde> o<verde>} \\
\\ 
\texttt{Fornisci la tua risposta solo in questo formato.} \\
\texttt{spiegazione: dettagli} \\
\texttt{tentativo: parola} \\
            }
        }
    }
     \\ \\

    \theutterance \stepcounter{utterance}  
    & \multicolumn{4}{p{0.6\linewidth}}{
        \cellcolor[rgb]{0.561,0.737,0.561}{
            \makecell[{{p{\linewidth}}}]{
                \texttt{\tiny{[P1$\rangle$GM]}}
                \texttt{spiegazione: le posizioni 2{-}5 sono o,n,n,o. La s è errata. "nonno" inizia con n e mantiene o,n,n,o.} \\
\texttt{tentativo: nonno} \\
            }
        }
    }
    & & \\ \\

    \theutterance \stepcounter{utterance}  
    & & & \multicolumn{2}{p{0.3\linewidth}}{
        \cellcolor[rgb]{0.9,0.9,0.9}{
            \makecell[{{p{\linewidth}}}]{
                \texttt{\tiny{[GM$|$GM]}}
                \texttt{attempts = 4} \\
\texttt{target\_word = nonno} \\
\texttt{guess = nonno} \\
\texttt{guess\_feedback = n<green> o<green> n<green> n<green> o<green>} \\
            }
        }
    }
    & & \\ \\

    \theutterance \stepcounter{utterance}  
    & & & \multicolumn{2}{p{0.3\linewidth}}{
        \cellcolor[rgb]{0.9,0.9,0.9}{
            \makecell[{{p{\linewidth}}}]{
                \texttt{\tiny{[GM$|$GM]}}
                \texttt{game\_result = WIN} \\
            }
        }
    }
    & & \\ \\

\end{supertabular}
}

\end{document}
